\documentclass{situationist-preprint}

\title{The Bank Cases}
\msid{EPI-13-b}
\msfamily{EPI-13}
\msdiscipline{Epistemology}
\mstype{Companion manuscript}
\mskeywords{contextualism; bank cases; stakes; safety; closeness; subject-sensitive invariantism}

\begin{document}
\maketitle

\begin{abstract}
DeRose's bank cases are the standard evidence that \enquote{know} is context-sensitive: the same speaker, with the same evidence that the bank is open on Saturday, truly says \enquote{I know} when little hangs on it and truly says \enquote{I don't know} when much does. We show that the account of EPI-13-a, on which context fixes which worlds are close, explains the cases, and that it explains them in a limited way. High stakes make close the worlds in which the bank has recently changed its hours, and nothing more remote. The sceptic gains nothing from the bank cases, and third-person variants favour contextualism over subject-sensitive invariantism.
\end{abstract}

\section{The cases}

In the low-stakes case, Keith and his wife drive past the bank on Friday, see long queues, and decide to return on Saturday; Keith says he knows the bank is open on Saturdays, having been there two weeks before. In the high-stakes case, a large cheque must be deposited before Monday, and his wife points out that banks do change their hours; Keith now says he does not know that the bank will be open \citep{DeRose1992}. Both utterances seem true, although Keith's evidence and belief are the same.

EPI-13-a locates the context-sensitivity of \enquote{know} in contextual closeness (Definition~2.1 there): the subject knows$_c$ $p$ if her belief is justified, true, and safe across the worlds close to hers by the standards of $c$ (Definition~2.3). Attending to a possibility that must be ruled out makes the worlds in which it obtains close (Assumption~2.9). The bank cases are an instance. Let $B$ be the proposition that the bank is open on Saturday, and $X_B$ the possibility that it has changed its hours since Keith's visit.

\begin{proposition}[Limited widening]\label{p:limited}
In the high-stakes context $c^{+}$, worlds in which $X_B$ holds and Keith still believes $B$ on his memory of the earlier visit are close, so Keith does not know$_{c^{+}}$ $B$. In the low-stakes context $c^{-}$, which attends to no such possibility, he knows$_{c^{-}}$ $B$, provided the bank has not in fact changed its hours. In both contexts he knows that he has hands.
\end{proposition}

\begin{proof}[Argument]
In $c^{+}$ the possibility $X_B$ is attended to as one that must be ruled out, since the cost of error makes it relevant to the decision under discussion. By the rule of attention, the $X_B$-worlds in which Keith believes $B$ on the same basis are close; $B$ is false there, so the belief is unsafe in $c^{+}$. In $c^{-}$ the close worlds are those in which the bank's routine is as Keith found it, and his belief is safe; it is justified and true, so he knows. No sceptical hypothesis about his hands is attended to in either context, so by Theorem~3.5 of EPI-13-a he knows$_c$ that he has hands in both.
\end{proof}

The widening is limited because high stakes do not attend to every possibility of error, only to those that bear on the practical question and are of a kind that ordinarily occurs: banks change their hours. Stakes alone do not bring the brain-in-a-vat world into view. This is why the bank cases do not support scepticism and why, by Lemma~3.2 of EPI-13-a, Keith's other beliefs deduced from his justified safe beliefs remain knowledge in each context.

\section{Objections}

\begin{objection}
Subject-sensitive invariantism explains the cases equally well: the subject's own stakes, not the attributor's context, raise the standard \citep{Stanley2005}.
\end{objection}

\begin{reply}
In the third-person variant, Keith's stakes are low, but his wife, whose cheque is due, says \enquote{Keith doesn't know that the bank will be open}; this seems true, though the subject's stakes have not changed \citep{DeRose1992}. Subject-sensitive invariantism must call it false. On the present view the attributor's context makes $X_B$-worlds close, and the verdict follows.
\end{reply}

\begin{objection}
Keith in the high-stakes case merely lacks the confidence needed to act, not knowledge.
\end{objection}

\begin{reply}
This predicts that \enquote{I know, but I had better check} should sound correct in the high-stakes case; it does not. The practical difference is mediated by what is attended to, and attention, through closeness, bears on knowledge \citep{Cohen1999}.
\end{reply}

\section{Conclusion}

The bank cases confirm that the standards for \enquote{know} shift with context, and show how far they shift: stakes widen the close worlds to include ordinary nearby error, and leave sceptical worlds remote unless they are themselves attended to.

\begin{thebibliography}{99}\small
\bibitem[Cohen(1999)]{Cohen1999} Cohen, S. (1999). Contextualism, skepticism, and the structure of reasons. \emph{Philosophical Perspectives}, 13, 57--89.
\bibitem[DeRose(1992)]{DeRose1992} DeRose, K. (1992). Contextualism and knowledge attributions. \emph{Philosophy and Phenomenological Research}, 52(4), 913--929.
\bibitem[Stanley(2005)]{Stanley2005} Stanley, J. (2005). \emph{Knowledge and Practical Interests}. Oxford University Press.
\end{thebibliography}

\end{document}
